\lectures{1,2,3}
\title{Random Variables}

\begin{document}

\subsection{Preliminaries}

First, getting some jargon and preliminary facts out of the way.
\begin{itemize}
    \item When we write $A \sim B$, we mean $A$ is a random variable
    distributed as some probability measure $B$.
    \begin{itemize}
        \item So we can write $X_1, X_2, \dots, X_n \sim \mathcal{N}(0, 1)$,
        but not $X \sim Y$ where $Y$ is a random variable.
        \item We often say samples $X_1, X_2, \dots, X_n \sim \mathbb{P}$ are
        independent, identically distributed, or i.i.d. for short.
        \item As shorthand for ``$A$ and $B$ are independent'',
        we may also write ``$A \perp B$''.
    \end{itemize}
    \item Given a random variable $X$\dots
    \begin{itemize}
        \item The sample mean is $\bar{X}$, whereas
        the true (population) mean is $\mu = \mathbb{E}[X]$.
        \begin{itemize}
            \item The sample mean $\bar{X}_n := \frac{1}{n}\sum_{i = 1}^n X_i$
            (assuming i.i.d. $X_1, X_2, \dots, X_n \sim \mathbb{P}$)
            is itself a random variable!
        \end{itemize}
        \item The sample variance is $s^2$, whereas
        the true (population) variance is $\sigma^2 = \mathbb{V}[X]$.
    \end{itemize}
    \item Some discrete distributions to know\dots
    \begin{itemize}
        \item $X \sim \text{Ber}(p)$ is a weighted coin flip,
        with $\Pr(X = 1) = p$ and $\Pr(X = 0) = 1 - p$.
        \item $X \sim \text{Bin}(n, p)$ is the sum of $n$ independent
        trials of $\text{Ber}(p)$.
        \item $X \sim \text{Pois}(\lambda)$ is given by
        $\Pr(X = k) = \frac{\lambda^k e^{-\lambda}}{k!}$.
    \end{itemize}
    \item Some continuous distributions to know\dots
    \begin{itemize}
        \item $X \sim \mathcal{U}(a, b)$ is the uniform distribution:
        all values in $[a, b]$ are equally likely.
        \item $X \sim \text{Exp}(\lambda)$ is the exponential distribution, with PDF
        $f(x) = \lambda e^{- \lambda x}$.
        \item $X \sim \mathcal{N}(\mu, \sigma^2)$ is the normal distribution, with PDF
        $f(x) = \frac{1}{\sigma \sqrt{2\pi}} \exp\left(-\frac{(x - \mu)^2}{2\sigma^2}\right)$.
    \end{itemize}
\end{itemize}

\begin{theorem}{Linearity}
    For any random variables $A$ and $B$, we have
    $\mathbb{E}[A + B] = \mathbb{E}[A] + \mathbb{E}[B]$.
    Furthermore, if $A$ and $B$ are independent, then
    $\mathbb{V}[A + B] = \mathbb{V}[A] + \mathbb{V}[B]$.
\end{theorem}

\begin{remark}
    If $A$ and $B$ are not independent, a counterexample to the latter
    of the two statements is $A = B$, say.
\end{remark}

An immediate corollary of the above fact is the following.

\begin{theorem}{Sample Mean Stats}
    Suppose $X_1, X_2, \dots, X_n \sim X$ are sampled i.i.d.,
    yielding a sample mean $\bar{X}_n$.  Suppose further
    that $\mathbb{E}[X_i] = \mu$ and $\mathbb{V}[X_i] = \sigma^2$.  Then
    $\mathbb{E}[\bar{X}_n] = \mu$ and
    $\mathbb{V}[\bar{X}_n] = \sigma^2 / n$.
\end{theorem}

\subsection{Convergence of Random Variables}

\textbf{Definition (Convergence in Probability).}
Given random variables $\{X_n\}_{n = 1}^{\infty}$
and a random variable $X$, we say ``$X_n$ converges to $X$ in probability''
(denoted $X_n \convprob X$) if:
\begin{quote}
    For every $\epsilon > 0$, we have $\Pr(|X_n - X| > \epsilon) \to 0$ as $n \to \infty$.
\end{quote}

\textbf{Definition (Convergence in Distribution).}
Given random variables $\{X_n\}_{n = 1}^{\infty}$ and a random variable $X$,
we say ``$X_n$ converges to $X$ in distribution'' (denoted $X_n \rightsquigarrow X$) if:
\begin{quote}
    For all $x$ where $x \mapsto \Pr(X \leq x)$ is continuous,
    we have $\Pr(X_n \leq x) \to \Pr(X \leq x)$ as $n \to \infty$.
\end{quote}

\begin{remark}
    Think of ``convergence in distribution'' as saying that the CDFs of $X_n$ and $X$
    get sufficiently close together.
\end{remark}

\textbf{Example.} Consider $X_n \sim \text{Ber}(\frac{1}{2})$
and $X \sim \text{Ber}(\frac{1}{2})$.  Then $X_n \rightsquigarrow X$ is true,
but $X_n \convprob X$ is not.

Expectedly, the latter of these is much less strict than the former.
Here's a not-completely-obvious theorem.

\begin{theorem}{Relationship Between Convergence}
    If $X_n \convprob X$, then $X_n \rightsquigarrow X$.
\end{theorem}

\begin{remark}
    Convergence in probability is \emph{possible}\dots
    but it's really uncommon.
    \begin{itemize}
        \item \textbf{Example.} Consider
        $X_1, X_2, \dots \overset{\text{i.i.d}}{\sim} \mathcal{U}(0, 1)$,
        and define $M_n = \max(X_1, \dots, X_n)$.
        Then $M_n \convprob 1$.
        \item \textbf{Example.} Consider $U \sim \mathcal{U}(0, 1)$,
        and define $X_n = U + \frac{1}{n}$.  Then $X_n \convprob U$.
    \end{itemize}
    These examples feel kind of hacky, though---and it turns out, they have to be!
    \begin{quote}
        \textbf{Fact.} If $X_n \perp X$ for all $n$ and $X_n \convprob X$,
        then $\Pr(X = c) = 1$ for some constant $c$.
    \end{quote}
    This isn't explicitly in the 18.650 curriculum, but we'll prove it later anyway,
    once we learn what we need to do so.
\end{remark}

Here's a completely-obvious theorem.

\begin{theorem}{Convergence to Constant}
    We have $X_n \convprob c$ if and only if $X_n \rightsquigarrow c$.
\end{theorem}

\subsection{LLN and CLT}

\begin{theorem}{Weak Law of Large Numbers}
    Consider i.i.d. random variables $X_1, X_2, \dots$ satisfying
    $\mathbb{E}[X_i] = \mu$.  Then their sample mean $\bar{X}_n$ converges
    in probability to $\mu$; that is, $\bar{X}_n \convprob \mu$.
\end{theorem}

\begin{proof}
    Chebyshev's; the point is that $\mathbb{V}[\bar{X}_n] \to 0$
    as $n \to \infty$. ~ $\blacksquare$
\end{proof}

\begin{remark}
    The ``Strong Law of Large Numbers'' roughly says, more strongly,
    that $\Pr(\lim_{n \to \infty} \bar{X}_n = \mu) = 1$.
\end{remark}

\begin{theorem}{Central Limit Theorem}
    Given $\{X_i\}_{i = 1}^n$ i.i.d. with mean $\mu$ and variance $\sigma^2$, we have:
    \[ \dfrac{\sqrt{n}}{\sigma}(\bar{X}_n - \mu) \rightsquigarrow \mathcal{N}(0, 1), ~
    \text{ or equivalently (hand-waving), } ~
     \bar{X}_n \approx \mathcal{N}\left(\mu, \dfrac{\sigma^2}{n}\right). \]
\end{theorem}

\begin{remark}
    We use some shorthand here: $X \rightsquigarrow \mathbb{P}$ is short
    for $X \rightsquigarrow Y \sim \mathbb{P}$.
\end{remark}

You can at least sanity-check CLT; the $\mathbb{E}$ and $\mathbb{V}$ clearly match.
But the proof is very nontrivial.

\subsection{Properties of Convergence}

The former of these two bullets is not hard to prove; the latter is pretty hard.

\begin{theorem}{Convergence of Sums and Products}
    Consider $\{X_n\}_{n = 1}^{\infty}$, $\{Y_n\}_{n = 1}^{\infty}$,
    $X$, and $Y$ random variables.
    \begin{itemize}
        \item If $X_n \convprob X$ and
        $Y_n \convprob Y$, then $X_n + Y_n \convprob X + Y$
        and $X_nY_n \convprob XY$.
        \item \emph{(Slutsky's Theorem)} If $X_n \rightsquigarrow X$ and
        $Y_n \convprob c$, then $X_n + Y_n \rightsquigarrow X + c$,
        and $X_n Y_n \rightsquigarrow Xc$.
    \end{itemize}
\end{theorem}

\begin{remark}
    Consider the following false statement.
    \begin{quote}
        If $X_n \rightsquigarrow X$ and $Y_n \rightsquigarrow Y$,
        then $X_n + Y_n \rightsquigarrow X + Y$.
    \end{quote}
    A counterexample is $X_n, X, Y \sim \mathcal{N}(0, 1)$ and $Y_n = -X_n$.
\end{remark}

To prove things about convergence, it helps to know the following fact.

\begin{theorem}{Continuous Mapping Theorem}
    If $X_n \convprob X$, then $g(X_n) \convprob g(X)$ for any continuous function $g$.
    Similarly, if $X_n \rightsquigarrow X$, then $g(X_n) \rightsquigarrow X$.
\end{theorem}

This lets us justify the so-called ``Delta method''.

\begin{theorem}{Delta Method}
    Consider a sequence of random variables $Y_n$ such that
    $\frac{\sqrt{n}}{\sigma}(Y_n - \mu) \rightsquigarrow Y \sim \mathcal{N}(0, 1)$.
    Then for any differentiable $g$ such that $g'(\mu) \neq 0$, we have:
    \[ \dfrac{\sqrt{n}}{\sigma} \left(g(Y_n) - g(\mu)\right) \rightsquigarrow \mathcal{N}(0, g'(\mu)^2). \]
    (Naturally, you should think to use the Delta method for $Y_n = \bar{X}_n$.)
\end{theorem}

More intuitively, the Delta method is just the CLT with a $g$ attached;
when $g(x) = x$, the Delta method is the CLT verbatim.

\begin{proof}
    Just Taylor expand $g$ around $\mu$.
    \[ \dfrac{\sqrt{n}}{\sigma} \left(g(Y_n) - g(\mu)\right)
    = g'(\mu)\left[\dfrac{\sqrt{n}}{\sigma}(Y_n - \mu)\right] +
    \dfrac{\sqrt{n}}{\sigma}R_n. \]
    The first term, by CLT, converges ($\rightsquigarrow$) to $\mathcal{N}(0, g'(\mu)^2)$.
    The second term converges ($\convprob$) to $0$. ~ $\blacksquare$
\end{proof}

\subsection{Applying Slutsky's Theorem}

Here's the set-up.

\begin{quote}
    \textbf{Set-Up.} Someone claims Harvard grades have a mean salary of below $120$K.
    To test this, you collect $X_1, X_2, \dots, X_{100}$ salaries (i.i.d.)
    and find that $\bar{X}_n = 121$K and $\hat{\sigma} = 0.3$K (in thousands).

    You'd like to say $\bar{X}_n \rightsquigarrow \mathcal{N}(120, 0.3^2/100)$;
    this would let you compute $\Pr(\bar{X}_n \geq 121)$ using the normal distribution.

    But the CLT doesn't immediately apply here---the \emph{true} standard
    deviation isn't known to be $\hat{\sigma}$.  What now?
\end{quote}

The point is that Slutsky's Theorem lets us say
$\bar{X}_n \rightsquigarrow \mathcal{N}(120, 0.3^2/100)$
anyway, even though we don't know the exact value of $\sigma$.

\begin{theorem}{Sample Variance is Unbiased}
    We have $\mathbb{E}[s^2] = \sigma^2$, where
    $s^2 := \frac{1}{n - 1} \sum_{i = 1}^n (X_i - \bar{X}_n)^2$
    is the sample variance.
\end{theorem}

\begin{proof}
    Just expand everything and apply linearity.  The details look like this:
    \begin{align*}
        \mathbb{E}\left[\sum_{i = 1}^n (X_i - \bar{X}_n)^2\right]
        & = \mathbb{E}\left[ \sum_{i = 1}^n (X_i - \mu)^2 \right]
        - \mathbb{E}\left[2(\bar{X}_n - \mu)\sum_{i = 1}^n (X_i - \mu)
        + n(\bar{X}_n - \mu)^2\right] \\
        & = \underbrace{\mathbb{E}\left[ \sum_{i = 1}^n (X_i - \mu)^2 \right]}_{n \sigma^2}
        - \underbrace{\mathbb{E}\left[n(\bar{X}_n - \mu)^2\right]}_{n\left(\frac{\sigma^2}{n}\right)}
        = (n - 1)\sigma^2.
    \end{align*}
    And that's why the sample variance has a $\frac{1}{n - 1}$ factor;
    if it were $\frac{1}{n}$, it would no longer be unbiased! ~ $\blacksquare$
\end{proof}

With the above, we can justify $\bar{X}_n \rightsquigarrow \mathcal{N}(120, 0.3^2/100)$ as follows:
\begin{enumerate}
    \item By the Law of Large Numbers, $\mathbb{E}[s^2] = \sigma^2$
    gives us $s^2 \convprob \sigma^2$.  (Notational issues aside\dots)
    \item Then the Continuous Mapping Theorem (with $g(x) = \sqrt{x}$)
    gives $\hat{\sigma} \convprob \sigma$.  In other words,
    $\frac{\sigma}{\hat{\sigma}} \convprob 1$.
    \item Also, the Central Limit Theorem tells us
    $\frac{\sqrt{n}}{\sigma} (\bar{X}_n - \mu) \rightsquigarrow \mathcal{N}(0, 1)$.
    \item Applying Slutsky's Theorem on points 2 and 3 yields
    $\frac{\sqrt{n}}{\hat{\sigma}} (\bar{X}_n - \mu) \rightsquigarrow \mathcal{N}(0, 1)$, done!
\end{enumerate}

\subsection{Convergence in Probability is Hard}

To close things off, let's prove that not-in-18.650-curriculum fact from earlier.

\begin{theorem}{Convergence in Probability is Hard}
    If $X_n \perp X$ for all $n$ and $X_n \convprob X$,
    then $\Pr(X = c) = 1$ for some constant $c$.
\end{theorem}

\begin{proof}
    We may assume $X_n$ and $X$ are bounded; if they're not,
    just replace them with $\tan^{-1}(X_n)$ and $\tan^{-1}(X)$,
    legal by the Continuous Mapping Theorem.  Then the independence yields:
    \[ \mathbb{E}[X^2] \leftarrow \mathbb{E}[X_n \cdot X] =
    \mathbb{E}[X_n] \cdot \mathbb{E}[X] \rightarrow \mathbb{E}[X]^2. \]
    Letting $n \to \infty$, we get $\mathbb{E}[X^2] = \mathbb{E}[X]^2$,
    or $\mathbb{V}[X] = 0$, which implies $X$ is almost surely constant. ~ $\blacksquare$
\end{proof}

\begin{remark}
    The convergence in probability---not true in convergence
    in distribution---comes from the step $X_n \cdot X \convprob X^2$.
\end{remark}

\end{document}